\documentclass[a4paper,12pt]{article}
\usepackage{graphicx} 
\begin{document}

\title{\vspace{-70pt}\begingroup  
    \centering
    \includegraphics[width=0.6\textwidth]{Imaxes/Simbolo_logo_cor transp.png}\\
    \large Grado en Intelixencia Artificial\\
    \large Sistemas Multiaxente\\[0.5em]
    \large \textbf{Práctica 2: Sistemas Multiagente}\par 
\endgroup}
\author{Nombre Apellido Apellido2 \\ Nombre Apellido Apellido2 \\ Grupo de clase (martes o jueves)}

\date{}
\maketitle
\vspace{-30pt}


\section*{\large Mundo Virtual}

Describid los aspectos más relevantes del mundo virtual que habéis desarrollado. Indicad las decisiones de diseño tomadas en relación con la distribución de habitaciones, pasillos, puertas u otros elementos clave del entorno. En caso de haber incorporado nuevos agentes, además de los cuatro policías requeridos, detallad su función y características en esta sección. La interacción de estos elementos con el sonido y la visión deberían describirse en esta sección.

\section*{\large Arquitectura individual}

Justificad brevemente la arquitectura adoptada para cada uno de los agentes del sistema. Explicad por qué se ha elegido dicha arquitectura, cómo está estructurada y cuál es su modo de funcionamiento. En caso de utilizar arquitecturas diferentes para distintos agentes, destacad las particularidades de cada una. Esto incluye como se incorporan las comunicaciones al comportamiento de los agentes.

\section*{\large Comunicación entre agentes}

Detallad la estructura de los mensajes intercambiados entre los agentes, así como el flujo de comunicación: cómo se crean, envían y reciben los mensajes. Indicad qué protocolos de comunicación habéis implementado. Si habéis realizado modificaciones sobre protocolos estándar, incluid un diagrama AUML que represente vuestro protocolo personalizado.

\section*{\large Formato del contenido de los lenguajes}

Especificad el lenguaje utilizado para representar el contenido de los mensajes intercambiados entre agentes. Describid cómo se estructura dicho contenido y qué mecanismos utilizáis para procesarlo y extraer la información relevante en el contexto del sistema.

\section*{\large Estrategia Multiagente}

Explicad en qué consiste la estrategia cooperativa implementada en vuestro sistema multiagente. Describid cómo interactúan los agentes para lograr un objetivo común, qué mecanismos de coordinación emplean y cómo se toman las decisiones a nivel global. Si existen roles (dinámicos) diferenciados, planes compartidos o estrategias adaptativas.

\end{document}
